\begin{definition}[Рёберная двусвязность]
Вершины \(u,v\) рёберно двусвязны, \(u\sim_E v\), если существуют
два \(uv\)-пути без общих рёбер. Для \(u=v\) допускаются два нулевых пути.
Пути можно считать простыми, удалив замкнутые участки.
\end{definition}
\begin{theorem}
Рёберная двусвязность — отношение эквивалентности на вершинах.
Его классы называются компонентами рёберной двусвязности.
\end{theorem}
\begin{proof}
Рефлексивность следует из нулевых путей, симметричность — из обращения.
Пусть \(u\sim_E v\), \(v\sim_E w\). Возьмём непересекающиеся по рёбрам
простые пути \(P_1,P_2\) из \(u\) в \(v\) и такие же пути \(Q_1,Q_2\)
из \(w\) в \(v\). Пусть \(a,b\) — первые вершины встречи \(Q_1,Q_2\)
с \(P_1\cup P_2\). До них пути не используют рёбер \(P_1,P_2\).
Если \(w\) уже лежит на объединении, допускаются нулевые начальные участки.

Если \(a,b\) можно отнести к разным \(P_i\), склеим начальный участок
каждого \(Q_i\) с участком соответствующего \(P_i\) до \(u\).
Получены два пути без общих рёбер. Иначе обе точки лежат только
на одном пути, пусть \(P_1\), в порядке \(u,\ldots,a,\ldots,b,\ldots,v\)
(при необходимости поменяем \(Q_1,Q_2\) местами).
Первый путь идёт из \(w\) по \(Q_1\) до \(a\), затем по \(P_1\) к \(u\).
Второй — по \(Q_2\) до \(b\), по \(P_1\) от \(b\) к \(v\), затем
по \(P_2\) к \(u\). Участки \(P_1\) не имеют общих рёбер, как и
начальные участки \(Q_1,Q_2\). При \(a=b\) рассуждение также работает.
Внутри отдельного пути возможные повторения удаляются.
Таким образом, \(u\sim_E w\).
\end{proof}
\begin{graphillustration}{Два случая в доказательстве транзитивности теоремы 4.1. Жирная сплошная и штриховая линии выделяют полученные $uw$-пути $R_1,R_2$ без общих рёбер; тонкие серые линии — неиспользуемые участки. Показаны только начальные участки $Q_1,Q_2$ до первых встреч с $P_1\cup P_2$.}{fig:04-transitivity}
\begin{tikzpicture}[font=\small]
\begin{scope}
  \node[gi note] at (2,2) {1. Точки встречи на разных путях};
  \node[gi vertex] (u) at (0,0) {$u$};
  \node[gi vertex] (v) at (4,0) {$v$};
  \node[gi vertex] (a) at (2,1.1) {$a$};
  \node[gi vertex] (b) at (2,-1.1) {$b$};
  \node[gi vertex] (w) at (2,0) {$w$};
  \draw[gi edge,black!45] (a)--(v);
  \draw[gi edge,black!45] (b)--(v);
  \draw[gi emphasis] (w)--(a)--(u);
  \draw[gi alternate] (w)--(b)--(u);
  \node[gi note] at (3.25,.9) {$P_1$};
  \node[gi note] at (3.25,-.9) {$P_2$};
  \node[gi note] at (2.45,.55) {$Q_1$};
  \node[gi note] at (2.45,-.55) {$Q_2$};
  \node[gi note] at (2,-2.9) {$R_1:w\xrightarrow{Q_1}a\xrightarrow{P_1}u$};
  \node[gi note] at (2,-3.55) {$R_2:w\xrightarrow{Q_2}b\xrightarrow{P_2}u$};
\end{scope}
\begin{scope}[shift={(7.2,0)}]
  \node[gi note] at (2,2) {2. Обе точки встречи на $P_1$};
  \node[gi vertex] (u) at (0,0) {$u$};
  \node[gi vertex] (v) at (4,0) {$v$};
  \node[gi vertex] (a) at (1.3,-.8) {$a$};
  \node[gi vertex] (b) at (2.7,-.8) {$b$};
  \node[gi vertex] (w) at (2,-2.1) {$w$};
  \draw[gi edge,black!45] (a)--(b);
  \draw[gi emphasis] (w)--(a)--(u);
  \draw[gi alternate] (w)--(b)--(v);
  \draw[gi alternate] (v) to[bend right=35] (u);
  \node[gi note] at (2,1.35) {$P_2$};
  \node[gi note] at (2,-.4) {$P_1$};
  \node[gi note] at (1.15,-1.6) {$Q_1$};
  \node[gi note] at (2.85,-1.6) {$Q_2$};
  \node[gi note] at (2,-2.9) {$R_1:w\xrightarrow{Q_1}a\xrightarrow{P_1}u$};
  \node[gi note] at (2,-3.55) {$R_2:w\xrightarrow{Q_2}b\xrightarrow{P_1}v\xrightarrow{P_2}u$};
\end{scope}
\end{tikzpicture}
\end{graphillustration}


\begin{definition}[Мост]
Ребро \(ab\) называется мостом, если его концы лежат в разных компонентах
рёберной двусвязности.
\end{definition}
\begin{lemma}[Мост, цикл и удаление]\label{lem:bridge}
Для ребра \(e=ab\) эквивалентны: оно является мостом; не лежит
ни на одном простом цикле; его удаление увеличивает число компонент связности.
\end{lemma}
\begin{proof}
Если \(e\) лежит на простом цикле, ребро и остальная дуга цикла дают
два пути между концами без общих рёбер, и удаление \(e\) не нарушает
связности: любое его использование заменяется обходом по циклу.
Если \(a\sim_E b\), хотя бы один из двух простых путей не использует
\(e\); вместе с \(e\) он образует простой цикл.
Наконец, если \(G-e\) всё ещё соединяет \(a,b\), простой путь между
ними вместе с \(e\) образует цикл. Если не соединяет, исходная компонента
распадается на две: любой простой путь от её вершины к \(a\) либо не использует
\(e\), либо до его прохождения соединяет вершину с \(b\).
\end{proof}
\begin{graphillustration}{Пример: удаление моста разделяет два треугольника. Пунктир отмечает удалённое ребро.}{fig:04-bridge}
\begin{tikzpicture}[font=\small]
\node[gi vertex] (a) at (0,0) {$a$};
\node[gi vertex] (p) at (-1,0.8) {$p$};
\node[gi vertex] (q) at (-1,-0.8) {$q$};
\node[gi vertex] (b) at (1.5,0) {$b$};
\node[gi vertex] (r) at (2.5,0.8) {$r$};
\node[gi vertex] (s) at (2.5,-0.8) {$s$};
\draw[gi edge] (a) -- (p);
\draw[gi edge] (p) -- (q);
\draw[gi edge] (q) -- (a);
\draw[gi edge] (b) -- (r);
\draw[gi edge] (r) -- (s);
\draw[gi edge] (s) -- (b);
\draw[gi emphasis] (a) -- (b);
\node[gi note] at (0.75,1.25) {$e=ab$ — мост};
\node[gi note] at (0.75,-1.45) {$c(G)=1$};
\node[gi note] at (4.3,0) {$\longrightarrow$};
\begin{scope}[shift={(6.5,0)}]
\node[gi vertex] (a) at (0,0) {$a$};
\node[gi vertex] (p) at (-1,0.8) {$p$};
\node[gi vertex] (q) at (-1,-0.8) {$q$};
\node[gi vertex] (b) at (1.5,0) {$b$};
\node[gi vertex] (r) at (2.5,0.8) {$r$};
\node[gi vertex] (s) at (2.5,-0.8) {$s$};
\draw[gi edge] (a) -- (p);
\draw[gi edge] (p) -- (q);
\draw[gi edge] (q) -- (a);
\draw[gi edge] (b) -- (r);
\draw[gi edge] (r) -- (s);
\draw[gi edge] (s) -- (b);
\draw[gi removed] (a) -- (b);
\node[gi note] at (0.75,-1.45) {$c(G-e)=2$};
\end{scope}
\end{tikzpicture}
\end{graphillustration}


\begin{theorem}[Удаление одного ребра и компоненты]
Вершины рёберно двусвязны тогда и только тогда, когда связаны и остаются
связанными после удаления любого одного ребра. Классы \(\sim_E\)
совпадают с компонентами графа после удаления всех мостов.
\end{theorem}
\begin{proof}
Два непересекающихся по рёбрам пути не разрушаются удалением одного ребра.
Если простой путь между \(u,v\) не проходит через мосты, концы каждого
его ребра рёберно двусвязны; по транзитивности \(u\sim_E v\).
Если же \(u,v\) находятся в разных компонентах после удаления мостов,
простой путь между ними проходит через некоторый мост.
Удаление этого моста разрывает связь между концами пути:
иначе его две стороны соединялись бы обходом. Это доказывает оба утверждения.
\end{proof}
